% Generated draft: requires mathematical and source review.
\documentclass[11pt]{article}
\usepackage[margin=25mm]{geometry}
\usepackage{fontspec}
\setmainfont{Latin Modern Roman}
\usepackage{amsmath,amssymb,mathtools,unicode-math}
\setmathfont{Latin Modern Math}
\usepackage{enumitem,needspace,xurl,hyperref}
\setlength{\parindent}{0pt}
\setlength{\parskip}{5pt}
\setlength{\emergencystretch}{3em}
\newcommand{\N}{\mathbb{N}}
\newcommand{\Z}{\mathbb{Z}}
\newcommand{\Q}{\mathbb{Q}}
\newcommand{\R}{\mathbb{R}}
\newcommand{\C}{\mathbb{C}}
\newcommand{\F}{\mathbb{F}}
\newcommand{\A}{\mathbb{A}}
\newcommand{\PP}{\mathbb{P}}
\newcommand{\E}{\mathbb{E}}
\newcommand{\eps}{\varepsilon}
\newcommand{\dd}{\,\mathrm{d}}
\begin{document}
\section*{UnsolvedMath ID: 6000017}
\textbf{Source problem code:} catalogue entry
\subsection{Definitions and mathematical statement}
A Hessian manifold is a manifold $M$ with a flat torsion-free affine connection $\nabla$ and a positive Riemannian metric $g$ locally expressible in affine coordinates as
\[
 g_{ij}=\frac{\partial^2\varphi}{\partial x_i\partial x_j}
\]
for local potentials $\varphi$. The source asks for compact examples. Its special classes ``hyperbolic'' and ``flat compact Hessian'' are not defined; since every Hessian connection is flat, the latter cannot be read as an additional connection-flatness condition alone. Nor is hyperbolic silently interpreted as constant negative Riemannian sectional curvature.

\textbf{Construction problem (Satoru Shimizu).} Construct compact Hessian manifolds other than direct products of hyperbolic and flat compact Hessian manifolds. In particular, use deformation theory of affine structures to construct one with no fiber-bundle structure.

The direct-product requirement concerns actual products in the intended geometric category. Being finitely covered by a product is a different condition and is not substituted. The excerpt does not specify whether equivalence is Hessian/affine isomorphism or merely diffeomorphism, nor its conventions concerning trivial factors.

The stronger nonfibering request needs a convention excluding trivial bundles over a point and the identity bundle. A natural nontrivial-bundle meaning is a smooth fiber bundle $M\to B$ with $0<\dim B<\dim M$ and positive-dimensional fibers, but the source may intend a narrower affine/geometric class. It does not limit the question to fibering over $S^1$. The requested deformation-theory method is retained, with preservation of a Hessian structure part of the desired construction, not an already-established deformation result.

The research summary claims recent finite-product-cover theorems in dimensions at most six and circle-fibration results for compact orientable Hessian and hyperbolic affine manifolds, mentioning periodic monodromy and Bieberbach alternatives. It gives no papers, theorem statements, or precise hypotheses. Those claims do not by themselves settle an actual-product question and are not accepted here as verified obstructions to the stronger construction. The summary says no explicit example was verified.
\subsection{Short English statement}
Construct compact Hessian manifolds beyond the specified direct-product families, especially a nonfibered example obtained by deforming affine structures. The meanings of hyperbolic/flat factors, product equivalence, and nontrivial fiber bundle are missing. Uncited finite-cover and circle-fibration claims do not authenticate a resolution.
\subsection{Sources}
\begin{itemize}
\item Local UnsolvedMath registry, permanent ID 6000017, source code catalogue entry, ``Nonproduct Compact Hessian Manifolds''; proposer: Satoru Shimizu, proposed year 1998. Status partially solved, classification PARTIAL-PROGRESS, updated 2026-08-06.
\item Furuhata, Matsuzoe, and Urakawa, \emph{Open Problems in Affine Differential Geometry and Related Topics} (1998), item 5(c), scanned PDF pages 1--2, \url{https://www.jstage.jst.go.jp/article/iis/4/2/4_2_125/_pdf/-char/en}. Reported manual OCR from page image and access: 2026-07-29. Current status NEEDS\_REVIEW; the article poses an open problem. The record says the J-STAGE article page marks it CC BY 4.0.
\item The summary's recent product-cover, circle-fibration, periodic-monodromy, and Bieberbach claims have no supplied references or exact hypotheses.
\end{itemize}
\end{document}
